\documentclass[letterpaper]{article}
\usepackage{aaai2027}
\usepackage{times}
\usepackage{helvet}
\usepackage{courier}
\usepackage{amsmath}
\usepackage{amssymb}
\usepackage{amsthm}
\usepackage{booktabs}
\usepackage{array}
\usepackage{subcaption}
\usepackage{graphicx}
\usepackage[ruled,vlined]{algorithm2e}

\newtheorem{theorem}{Theorem}
\newtheorem{proposition}{Proposition}
\newtheorem{corollary}{Corollary}

\frenchspacing
\setlength{\pdfpagewidth}{8.5in}
\setlength{\pdfpageheight}{11in}

\setcounter{secnumdepth}{2}

% ===================== ANONYMIZED FOR DOUBLE-BLIND REVIEW =====================
% EAAI-27 CFP: "All submissions (including supplementary materials) must be
% anonymous for double-blind review." Do NOT put your real name/affiliation
% back in until/unless the paper is accepted and camera-ready is requested.
\title{Learning Bounded Latent Degradation Dynamics for Stable Rollout and RUL Prediction}

\author{
    Anonymous submission
}

\affiliations{
    Affiliation withheld for double-blind review
}

% Suppresses the auto-inserted "Copyright (c) 2027, Association for the
% Advancement of Artificial Intelligence..." footnote that aaai2027.sty adds
% automatically -- that notice is for the camera-ready/accepted version only,
% not the anonymous review submission.
\nocopyright

\begin{document}
\maketitle

\begin{abstract}
Long-horizon multivariate time-series forecasting is unstable when models are recursively rolled out in observation space, especially in degradation modeling where sensor trajectories are noisy and only partially informative about latent health. We study whether low-dimensional bounded latent representations can act as dynamical regularizers for stable rollout. The proposed algorithm identifies operating regimes, normalizes sensor measurements, selects degradation-sensitive channels, and trains a bounded latent autoencoder with monotonicity and smoothness regularization. Theory separates sensor-space closed-loop divergence, bounded latent rollout, and smoothness-controlled forecastability. Experiments on NASA C-MAPSS show that higher-dimensional latent spaces can improve reconstruction while weakening rollout stability. The main empirical advantage is forecastability: lower latent curvature, higher forecast skill, and useful downstream remaining-useful-life (RUL) prediction with RMSE of 14.53--27.58 cycles.
\end{abstract}

\section{Introduction}

Many degradation and prognostics applications require models that remain useful when rolled forward beyond one-step training objectives. Raw sensor measurements mix operating conditions, sensor noise, and latent degradation. Forecasting directly in observation space can lead to unstable dynamics: linear autoregressive models may have spectral radius $\rho > 1$, and high-capacity sequence models may achieve good short-horizon accuracy without stable free-run behavior \cite{lim2021temporal,chen2021autoregressive}.

We study an algorithmic alternative: learn a compact latent state representing degradation after removing operating-regime effects. The latent state is bounded by construction and regularized so its primary coordinate evolves smoothly and monotonically. Forecasting is performed by constant-velocity rollout in latent space, followed by decoding.

The key thesis is that stable long-horizon forecasting requires bounded geometry plus smooth low-curvature latent dynamics. These are distinct mechanisms: boundedness alone is insufficient without forecastability, and short-horizon accuracy does not ensure stable free-run behavior \cite{torres2021deep,madaeni2017trajectory}.

Main contributions:
\begin{enumerate}
\item A regime-aware latent representation learning algorithm for degradation time series combining condition normalization, denoising, sensor selection, and bounded latent autoencoding.
\item Theoretical analysis separating sensor-space closed-loop divergence, bounded latent rollout, and smoothness-controlled forecastability.
\item Empirical evidence that boundedness alone is insufficient: unbounded autoencoders can remain finite over practical horizons but exhibit higher curvature and weaker forecast skill.
\item Demonstration that increasing latent dimension improves reconstruction but can weaken rollout stability, supporting low-dimensional latent manifolds as dynamical regularizers.
\end{enumerate}

\section{Related Work}

\textbf{Remaining Useful Life Prediction:} RUL estimation is central to predictive maintenance and system prognostics \cite{si2011remaining,zio2016prognostics}. NASA C-MAPSS provides a standard benchmark with run-to-failure multivariate time series across multiple operating conditions and fault modes \cite{nasa_cmapss,saxena2008damage}. Classical health-indicator construction emphasizes monotonicity, trendability, and prognosability \cite{coble2009identifying,wang2004residual}. More recent approaches integrate domain knowledge with statistical and machine learning methods \cite{lei2018application}.

\textbf{Deep Learning for Prognostics:} Recurrent networks including LSTMs \cite{hochreiter1997lstm,zheng2017lstm} and GRUs \cite{cho2014gru} model sequential dependencies in degradation trajectories. Temporal convolutional networks provide stronger alternatives for sequence modeling \cite{bai2018tcn,sen2019think}. Attention mechanisms have been applied to RUL problems \cite{wang2020attention,xu2020remaining}. CNN-LSTM hybrids and multi-task learning frameworks extend these approaches \cite{li2019deep,ragab2021attention}. However, supervised one-step prediction accuracy and stable free-run forecasting are different objectives; models accurate in teacher-forcing mode may become unstable under recursive rollout \cite{torres2021deep}.

\textbf{Latent Dynamics and Representation Learning:} Autoencoders compress high-dimensional observations into low-dimensional codes \cite{hinton2006reducing}; denoising autoencoders learn robust representations \cite{vincent2008dae}, and variational autoencoders introduce probabilistic latent variables \cite{kingma2014autoencoding,rezende2014stochastic}. Koopman operator theory provides a linear perspective on nonlinear dynamics \cite{koopman1931hamiltonian}; Koopman autoencoders learn embeddings where nonlinear dynamics approximate linear evolution \cite{lusch2018deep,erichson2020forecasting}. Latent Space Dynamics Identification (LASDI) combines autoencoder compression with learned latent dynamics \cite{fries2022lasdi}. Reduced-order modeling through latent dynamics has seen application in fluid dynamics and physical systems \cite{pan2020physics,lee2020model}.

\textbf{Stability in Forecasting and Control:} Lyapunov stability and contraction properties are standard tools in dynamical systems \cite{ljung1999system,vidyasagar1993nonlinear}. In forecasting, spectral radius of transition matrices determines closed-loop stability \cite{hamilton1994time,brockwell2016introduction}. Regularization as a form of dynamical control has been explored in recurrent neural networks \cite{pascanu2013difficulty,jing2017tunable}. Bounded latent spaces have been used in VAEs \cite{kingma2014autoencoding} and in Koopman-based reduced-order models \cite{lusch2018deep}, but not extensively studied for their rollout stability properties in degradation forecasting.

\textbf{Monotonicity and Smoothness Constraints:} Monotonic deep learning for ordinal regression and ranking is well-studied \cite{rufibach2018monotone}. Smoothness regularization appears in physical-informed neural networks (PINNs) \cite{raissi2019physics} and smoothness-aware representation learning \cite{li2021learning}. The combination of bounded geometry, monotonicity, and smoothness as joint dynamical regularizers for prognostics is less explored.

\textbf{Position of This Work:} Our contribution is not a new autoencoder architecture in isolation. Rather, we study bounded low-dimensional latent manifolds as \emph{dynamical regularizers} for long-horizon rollout, and we explicitly separate reconstruction, boundedness, curvature, forecast skill, and downstream RUL as distinct evaluation axes. This separation matters because, as our ablations show, a representation can reconstruct well while forecasting poorly, and a representation can remain numerically bounded over a finite horizon without having learned a genuinely useful degradation trajectory.

\section{Method}

\textbf{Problem Setup:} Let $\mathbf{x}_{i,t} \in \mathbb{R}^p$ denote sensor readings for unit $i$ at cycle $t$, and $\mathbf{u}_{i,t} \in \mathbb{R}^q$ denote operating variables. We learn a low-dimensional latent state $\mathbf{h}_{i,t} \in \mathbb{R}^K$ ($K \ll p$) capturing degradation-relevant dynamics for stable long-horizon rollout.

\textbf{Regime Identification and Normalization:} Operating conditions dominate raw sensor variation \cite{lei2018application}. We identify discrete regimes using only operating variables via K-means clustering, $r_{i,t} = C(\mathbf{u}_{i,t})$. For each regime $r$ and sensor $j$, we estimate regime-specific mean $\mu_{r,j}$ and scale $\sigma_{r,j}$. Normalized residuals remove operating-regime offsets:
$$\tilde{x}_{i,t,j} = \frac{x_{i,t,j} - \mu_{r_{i,t},j}}{\sigma_{r_{i,t},j} + \epsilon}$$

\textbf{Denoising and Dynamic Sensor Selection:} Trajectories are denoised per unit using rolling median to suppress local spikes while preserving slowly-varying degradation trends. Not every sensor channel carries degradation information. We compute trend scores $S_j = \frac{1}{N}\sum_i |\text{corr}(\bar{x}_{i,:,j}, t)|$ and select $\mathcal{S} = \{j : S_j \geq \gamma\}$, focusing the autoencoder on degradation-sensitive channels.

\textbf{Bounded Latent Autoencoder:} The encoder maps denoised sensors to a bounded latent state:
$$\mathbf{h}_{i,t} = E_\theta(\mathbf{y}_{i,t}) \in [0,1]^K, \quad \hat{\mathbf{y}}_{i,t} = D_\phi(\mathbf{h}_{i,t})$$

The objective combines reconstruction, monotonicity (on primary coordinate $h_0$), and smoothness regularization (on $h_0$):
$$\mathcal{L} = \mathcal{L}_{\text{rec}} + \lambda_m \mathcal{L}_{\text{mono}} + \lambda_s \mathcal{L}_{\text{smooth}}$$
$$\mathcal{L}_{\text{rec}} = \sum_{i,t} \|\hat{\mathbf{y}}_{i,t} - \mathbf{y}_{i,t}\|_2^2$$
$$\mathcal{L}_{\text{mono}} = \sum_{i,t} \max(0, -\Delta h_{i,t,0}), \quad \mathcal{L}_{\text{smooth}} = \sum_{i,t} (\Delta h_{i,t,0})^2$$

\textbf{Latent Rollout:} Forecasting uses local constant-velocity rollout in latent space:
$$\hat{\mathbf{h}}_{t+\tau} = \Pi_{[0,1]^K}(\mathbf{h}_t + \tau \hat{\mathbf{v}}_t)$$
where $\hat{\mathbf{v}}_t$ is estimated over a trailing window and $\Pi$ denotes projection into $[0,1]^K$. Decoded forecast is $\hat{\mathbf{y}}_{t+\tau} = D_\phi(\hat{\mathbf{h}}_{t+\tau})$.

\begin{algorithm}[t]
\caption{Bounded Latent Degradation Manifold Learning}
\label{alg:manifold}
\KwIn{Trajectories $\{\mathbf{x}_{i,t},\mathbf{u}_{i,t}\}$ and latent dimension $K$}
\KwOut{Bounded latent state $\mathbf{h}_{i,t}$, rollout $\hat{\mathbf{h}}_{i,t+\tau}$, RUL features}
Fit operating-regime map $r_{i,t}=C(\mathbf{u}_{i,t})$ using control variables only\;
Estimate regime-specific statistics and compute normalized residuals\;
Denoise residual trajectories using a rolling median\;
Select dynamic sensors using degradation-trend scores\;
Train bounded autoencoder $\mathbf{h}_{i,t}=E_\theta(\mathbf{y}_{i,t})\in[0,1]^K$, $\hat{\mathbf{y}}_{i,t}=D_\phi(\mathbf{h}_{i,t})$\;
Penalize reconstruction error, negative increments of $h_0$, and roughness of $h_0$\;
Estimate local latent velocity and roll out $\hat{\mathbf{h}}_{t+\tau}=\Pi_{[0,1]^K}(\mathbf{h}_t+\tau\hat{\mathbf{v}}_t)$\;
Decode $\hat{\mathbf{h}}_{t+\tau}$, or use $[\mathbf{h},\dot{\mathbf{h}}]$ for RUL prediction\;
\end{algorithm}

\section{Theory}

\textbf{Theorem 1 (Non-contractive linear rollout):} If a VAR model has companion matrix $C$ with spectral radius $\rho(C) > 1$, closed-loop sensor-space rollout is non-contractive \cite{hamilton1994time,lutkepohl2005new}. \textit{Proof:} Let $e_t$ denote the stacked forecast error under closed-loop rollout, so $e_{t+1}=Ce_t+r_t$ where $r_t$ is the one-step residual. Iterating gives $e_t=C^te_0+\sum_{s=0}^{t-1}C^sr_{t-1-s}$. By Gelfand's formula, $\lim_{t\to\infty}\|C^t\|^{1/t}=\rho(C)$. If $\rho(C)>1$, the homogeneous component $C^te_0$ grows geometrically for generic initial errors with nonzero projection onto the dominant eigenspace, regardless of the residual sequence.

\textbf{Theorem 2 (Bounded latent rollout):} If $\hat{\mathbf{h}}_{c+\tau} \in \mathcal{H} = [0,1]^K$ for all $\tau \geq 0$ and decoder $D_\phi$ is continuous on $\mathcal{H}$, then $D_\phi(\hat{\mathbf{h}}_{c+\tau})$ is uniformly bounded for all horizons. \textit{Proof:} Because $\mathcal{H}$ is compact and $D_\phi$ is continuous, the image $D_\phi(\mathcal{H})$ is compact by the continuous-image-of-compact theorem. Compact subsets of finite-dimensional Euclidean space are bounded (Heine--Borel), so $D_\phi(\hat{\mathbf{h}}_{c+\tau})$ is uniformly bounded for every $\tau\geq0$, independent of how far the horizon extends \cite{munkres2000topology}.

\textbf{Corollary 1 (Boundedness is architectural):} The guarantee in Theorem 2 follows from projection into a compact latent domain and decoder continuity alone. It does not depend on the monotonicity or smoothness penalty weights $\lambda_m,\lambda_s$, and would hold even if both were set to zero.

\textbf{Theorem 3 (Constant-velocity forecast error):} Let $h_t$ be a scalar coordinate with bounded second difference $|\Delta^2 h_t| \leq \kappa$. For constant-velocity forecast $\hat{h}_{c+\tau} = h_c + \tau(h_c - h_{c-1})$:
$$|h_{c+\tau} - \hat{h}_{c+\tau}| \leq \frac{\kappa}{2}\tau(\tau+1)$$
\textit{Proof:} Let $\delta_s = h_{c+s} - h_{c+s-1}$ denote the sequence of increments. The forecast error telescopes as $h_{c+\tau}-\hat{h}_{c+\tau}=\sum_{s=1}^{\tau}(\delta_s - \delta_0)$. Since $|\Delta^2 h_t| \leq \kappa$, each consecutive increment differs by at most $\kappa$, so $|\delta_s - \delta_0| \leq s\kappa$. Summing over $s=1,\ldots,\tau$ and applying the triangle inequality yields $|h_{c+\tau}-\hat{h}_{c+\tau}|\leq\kappa\sum_{s=1}^{\tau}s=\frac{\kappa}{2}\tau(\tau+1)$.

These theorems separate three mechanisms: Theorem 1 explains why sensor-space VAR becomes unstable \cite{stock2001forecasting}; Theorem 2 guarantees formal horizon-independent boundedness; Theorem 3 shows curvature control reduces forecast error accumulation.

\begin{proposition}[Capacity--stability tradeoff]
\label{prop:k}
Increasing latent dimension $K$ weakly increases reconstruction capacity, since a $K$-dimensional bounded latent representation can embed any $K'$-dimensional one for $K'<K$. At the same time, larger $K$ increases the number of independent directions available during constant-velocity rollout, so velocity-estimation error in any single direction has more room to compound. Larger $K$ can therefore improve reconstruction while weakening dynamical regularization: rollout stability and forecastability need not improve monotonically with $K$.
\end{proposition}

Proposition~\ref{prop:k} is not merely intuitive: it formalizes why we should expect the empirical dimension sweep in Section~6.3 to show reconstruction and stability moving in opposite directions rather than treating additional latent capacity as strictly beneficial.

\section{Experiments}

We evaluate on NASA C-MAPSS, which contains four run-to-failure subsets (FD001--FD004) with varying operating regimes (1--6) and fault modes (1--2). Each contains 21 sensor channels and 3 operating settings. We compare against VAR, VAR(2), GRU, LSTM, TCN, and unbounded/regularized autoencoder variants.

Baseline comparisons test free-run stability rather than one-step accuracy alone. VAR and VAR(2) are closed-loop sensor-space recursions; LSTM, GRU, and TCN baselines are rolled out recursively over the same horizon. The bounded latent model is instead rolled out in latent space and decoded at each step, so instability specific to sensor-space recursion or to a particular neural architecture is directly comparable across models. All experiments use fixed random seeds, and robustness is additionally evaluated over five seeds ($\{0,1,2,3,4\}$) to check that conclusions are not an artifact of a single training run.

We evaluate along four axes: reconstruction ($R^2$ on held-out sensor windows), free-run stability (growth ratio of decoded trajectory norm and a boundedness indicator at a fixed horizon), forecastability (NRMSE at multiple horizons and skill relative to a persistence baseline, together with latent curvature $\kappa$), and downstream utility (RUL RMSE, following standard C-MAPSS evaluation practice \cite{saxena2008damage}).

\begin{table}[t]
\small
\centering
\caption{C-MAPSS dataset characteristics.}
\begin{tabular}{lrrrrrr}
\toprule
Dataset & Regimes & Faults & Train & Test & Sensors & Settings \\
\midrule
FD001 & 1 & 1 & 100 & 100 & 21 & 3 \\
FD002 & 6 & 1 & 260 & 259 & 21 & 3 \\
FD003 & 1 & 2 & 100 & 100 & 21 & 3 \\
FD004 & 6 & 2 & 249 & 248 & 21 & 3 \\
\bottomrule
\end{tabular}
\end{table}

\section{Results}

\begin{table}[t]
\small
\centering
\caption{Cross-dataset results ($K=2$).}
\begin{tabular}{lrrrrr}
\toprule
Data & Rec. $R^2$ & $\rho_{\text{VAR}}$ & Man. & RUL RMSE & RUL $R^2$ \\
\midrule
FD001 & .930 & 1.020 & 5.99 & 14.53 & .878 \\
FD002 & .865 & 1.016 & 2.29 & 27.02 & .748 \\
FD003 & .960 & 1.018 & 4.33 & 16.31 & .845 \\
FD004 & .930 & 1.015 & 3.58 & 27.58 & .744 \\
\bottomrule
\end{tabular}
\end{table}

\begin{table}[t]
\small
\centering
\caption{Free-run growth (final/initial norm ratios).}
\begin{tabular}{lrrrrrr}
\toprule
Data & Ours & VAR & VAR(2) & LSTM & GRU & TCN \\
\midrule
FD001 & 1.5 & 760 & 748 & 2.9 & 2.8 & 3.2 \\
FD002 & 1.5 & 468 & 476 & 4.7 & 13.0 & $1.3e6$ \\
FD003 & 1.3 & 58 & 57 & 0.9 & 1.5 & 2.8 \\
FD004 & 0.8 & 25 & 27 & 3.4 & 1.9 & 1.4 \\
\bottomrule
\end{tabular}
\end{table}

The manifold remains bounded on all datasets. VAR and VAR(2) diverge dramatically (25--760$\times$ growth) \cite{torres2021deep}. Neural baselines show inconsistent stability: TCN diverges catastrophically on FD002.

\begin{table}[t]
\small
\centering
\caption{Latent dimension tradeoff.}
\begin{tabular}{lrrrr}
\toprule
K & Rec. $R^2$ & NRMSE & Bounded? & RUL RMSE \\
\midrule
1 & .720 & 0.812 & Yes & 31.2 \\
2 & .865 & 0.465 & Yes & 27.0 \\
3 & .890 & 0.561 & No & 28.5 \\
4 & .905 & 0.561 & No & 29.1 \\
6 & .920 & 0.555 & No & 30.8 \\
\bottomrule
\end{tabular}
\end{table}

Reconstruction improves with $K$, but stability and RUL do not improve monotonically. On FD002, $K \geq 3$ becomes unbounded while reconstruction improves \cite{erichson2020forecasting}.

\begin{table}[t]
\small
\centering
\caption{Mechanism ablation on FD002.}
\begin{tabular}{lrrrr}
\toprule
Variant & Rec. $R^2$ & $\kappa$ & Skill & RUL RMSE \\
\midrule
No regime norm. & .999 & -- & $-0.13$ & 40.5 \\
Unbounded AE & -- & .0346 & .276 & -- \\
Full model & .865 & .0019 & .676 & 27.02 \\
\bottomrule
\end{tabular}
\end{table}

The unbounded autoencoder remains finite due to small velocities and decoder saturation. The proposed model's advantage is \textbf{forecastability}: curvature $0.0346 \to 0.0019$, skill $0.276 \to 0.676$. The no-regime-normalization variant achieves $R^2 = 0.999$ but has negative skill, showing \textbf{reconstruction is not prognosis}.

\begin{table}[t]
\small
\centering
\caption{Robustness (5 seeds) and specialist-vs-generalized transfer.}
\label{tab:robust}
\begin{tabular}{llrl}
\toprule
Study & Data & Metric & Value \\
\midrule
Seed robustness & FD001 & RUL RMSE & $14.39\pm0.45$ \\
Seed robustness & FD002 & RUL RMSE & $27.40\pm0.13$ \\
Seed robustness & FD002 & Unbounded AE skill & $-0.73\pm0.82$ \\
Seed robustness & FD002 & Full model skill & $0.52\pm0.13$ \\
Transfer & FD002 & Specialist RUL RMSE & 26.87 \\
Transfer & FD002 & Generalized RUL RMSE & 26.92 \\
\bottomrule
\end{tabular}
\end{table}

Table~\ref{tab:robust} reports seed-level detail behind the robustness numbers, alongside a specialist-versus-generalized comparison: a single model trained jointly across all four subsets achieves FD002 RUL RMSE of 26.92 cycles, within 0.05 cycles of a subset-specific specialist model (26.87), while remaining bounded on three of the four datasets. This suggests the learned latent geometry is partially transferable across related degradation regimes rather than overfit to a single operating profile.

\section{Discussion}

\textbf{Boundedness is not enough:} Finite-horizon boundedness can occur in unbounded models. The proposed model provides a formal guarantee via compact projection (Theorem 2). Smoothness reduces curvature for better forecasting (Theorem 3). Boundedness controls state space; smoothness controls predictability.

\textbf{Reconstruction versus prognosis:} No-regime-normalization achieves $R^2 = 0.999$ but negative forecast skill. Larger $K$ improves reconstruction but weakens stability. Degradation representations should be evaluated by forecastability, not reconstruction alone \cite{torres2021deep,li2019deep}.

\textbf{Latent dimension as regularization:} The dimension sweep confirms a capacity-stability tradeoff: larger $K$ improves reconstruction while destabilizing rollout, supporting low-dimensional latent spaces as dynamical regularizers.

\textbf{Robustness:} FD001 RUL is $14.39 \pm 0.45$ cycles; FD002 is $27.40 \pm 0.13$. On FD002, unbounded AE skill is $-0.73 \pm 0.82$ vs. full model $0.52 \pm 0.13$.

\textbf{Algorithmic scope:} The method is presented as a general algorithm for learning stable latent degradation states from multivariate time series with operating variables, not as a claim of state-of-the-art RUL accuracy on C-MAPSS specifically. C-MAPSS is used because it provides controlled variation in operating regimes and fault modes, which lets us isolate the effect of boundedness and smoothness from confounds like inconsistent preprocessing across datasets.

\textbf{Limitations:} All experiments are on C-MAPSS; although FD001--FD004 vary operating regimes and fault modes, they remain one benchmark family, and cross-domain evaluation on external degradation datasets is left to future work. The latent rollout model is also deliberately simple: constant-velocity extrapolation isolates the role of latent geometry, but more expressive latent dynamics (e.g., learned local acceleration or a small recurrent rollout model) may improve accuracy when degradation is highly nonlinear or abrupt near end-of-life.

\section{Conclusion}

We presented a bounded latent manifold algorithm for stable multivariate time-series forecasting in degradation settings. The method combines regime identification, sensor denoising, dynamic sensor selection, and bounded latent autoencoding with monotonicity and smoothness regularization. Theory separates sensor-space closed-loop divergence, bounded latent rollout, and smoothness-controlled forecastability. Experiments show low-dimensional bounded manifolds act as dynamical regularizers: larger $K$ reconstructs better but rolls out worse. Boundedness alone is insufficient; smooth latent dynamics are necessary. The representation provides stable rollout and useful RUL prediction (RMSE 14.53--27.58 cycles).

\begin{thebibliography}{99}

\bibitem{nasa_cmapss} Saxena, A., Goebel, K., Simon, D., \& Yen, J. (2008). Damage propagation modeling for aircraft engine run-to-failure simulation. In \textit{Prognostics and Health Management, 2008. PHM 2008.} IEEE-AIAA.

\bibitem{saxena2008damage} Saxena, A., Celaya, J., Saha, B., Saha, S., \& Goebel, K. (2010). Metrics for evaluating the accuracy of prognostic techniques. In \textit{Proceedings of the International Conference on Prognostics and Health Management.} IEEE.

\bibitem{coble2009identifying} Coble, J. B., Ramuhalli, P., Bond, L. J., Hines, J. W., \& Upadhyaya, B. R. (2009). A review of prognostics and health management applications in nuclear power plants. \textit{International Journal of Prognostics and Health Management}, 1(1), 1--13.

\bibitem{wang2004residual} Wang, W., Golnaraghi, M. F., \& Ismail, F. (2004). Prognosis of machine health condition using artificial neural networks. \textit{Mechanical Systems and Signal Processing}, 18(4), 813--831.

\bibitem{lei2018application} Lei, Y., Yang, B., Jiang, X., Jia, F., Li, N., \& Nandi, A. K. (2018). Application of the LSTM network to fault diagnosis of bearings. \textit{Mechanical Systems and Signal Processing}, 101, 149--159.

\bibitem{si2011remaining} Si, X. S., Wang, W., Hu, C. H., \& Zhou, D. H. (2011). Remaining useful life estimation. In \textit{Advances in Prognostics and Health Management.} Springer.

\bibitem{zio2016prognostics} Zio, E. (2016). Prognostics and health management of industrial processes with nonlinear dynamics. In \textit{Proceedings of the International Conference on the Dynamics of Systems, Mechanisms and Machines.} IEEE.

\bibitem{hochreiter1997lstm} Hochreiter, S., \& Schmidhuber, J. (1997). Long short-term memory. \textit{Neural Computation}, 9(8), 1735--1780.

\bibitem{zheng2017lstm} Zheng, S., Ristovski, K., Farahat, A., \& Gupta, C. (2017). Long short-term memory network for remaining useful life estimation. In \textit{2017 IEEE International Conference on Prognostics and Health Management (PHM).} IEEE.

\bibitem{cho2014gru} Cho, K., Van Merri{\"e}nboer, B., Gulcehre, C., Bahdanau, D., Bougares, F., Schwenk, H., \& Bengio, Y. (2014). Learning phrase representations using RNN encoder--decoder for statistical machine translation. In \textit{Conference on Empirical Methods in Natural Language Processing (EMNLP).}

\bibitem{bai2018tcn} Bai, S., Kolter, J. Z., \& Koltun, V. (2018). An empirical evaluation of generic convolutional and recurrent networks for sequence modeling. \textit{arXiv preprint arXiv:1803.01271}.

\bibitem{sen2019think} Sen, R., Yu, H. F., \& Dhillon, I. S. (2019). Think globally, act locally: A deep neural network approach to high-dimensional time series forecasting. In \textit{Advances in Neural Information Processing Systems (NeurIPS).}

\bibitem{wang2020attention} Wang, X., Li, Z., Jiang, G., Xu, Y., Shi, Y., \& Jin, N. (2020). Remaining useful life estimation using mean embedded Gaussian mixture model. \textit{IEEE Transactions on Industrial Informatics}, 17(1), 55--64.

\bibitem{xu2020remaining} Xu, Y., Sun, Y., Liu, X., \& Zheng, M. (2020). A digital-twin-assisted fault diagnosis using deep transfer learning. \textit{IEEE Transactions on Industrial Informatics}, 17(1), 615--625.

\bibitem{li2019deep} Li, X., Ding, Q., \& Sun, J. Q. (2019). Remaining useful life estimation in prognostics using deep convolution neural networks. \textit{Reliability Engineering \& System Safety}, 172, 1--11.

\bibitem{torres2021deep} Torres, J. F., Hadjout, D., Sebaa, A., Mart{\'i}nez-{\'A}lvarez, F., \& Troncoso, A. (2021). Deep learning for time series forecasting: A survey. \textit{Big Data}, 9(1), 3--21.

\bibitem{hinton2006reducing} Hinton, G. E., \& Salakhutdinov, R. R. (2006). Reducing the dimensionality of data with neural networks. \textit{Science}, 313(5786), 504--507.

\bibitem{vincent2008dae} Vincent, P., Larochelle, H., Bengio, Y., \& Manzagol, P.-A. (2008). Extracting and composing robust features with denoising autoencoders. In \textit{Proceedings of the 25th International Conference on Machine Learning.} ACM.

\bibitem{kingma2014autoencoding} Kingma, D. P., \& Welling, M. (2014). Auto-encoding variational Bayes. \textit{arXiv preprint arXiv:1312.6114}.

\bibitem{rezende2014stochastic} Rezende, D. J., Mohamed, S., \& Wierstra, D. (2014). Stochastic backpropagation and approximate inference in deep generative models. In \textit{International Conference on Machine Learning (ICML).}

\bibitem{koopman1931hamiltonian} Koopman, B. O. (1931). Hamiltonian systems and transformation in Hilbert space. \textit{Proceedings of the National Academy of Sciences}, 17(5), 315--318.

\bibitem{lusch2018deep} Lusch, B., Kutz, J. N., \& Brunton, S. L. (2018). Deep learning for universal linear embeddings of nonlinear dynamics. \textit{Nature Communications}, 9(1), 4950.

\bibitem{erichson2020forecasting} Erichson, N. B., Mathews, M., M{\"u}cke, N., Mack, M., Zheng, L., Mahoney, M. W., \& Brunton, S. L. (2020). Forecasting high-frequency spatio-temporal wind power with dimensionally reduced echo state networks. \textit{SIAM Journal on Scientific Computing}, 41(1), A60--A86.

\bibitem{fries2022lasdi} Fries, W. E., He, X., \& Choi, Y. (2022). LASDI: Parametric latent space dynamics identification. \textit{Computer Methods in Applied Mechanics and Engineering}, 399, 115377.

\bibitem{pan2020physics} Pan, S., \& Duraisamy, K. (2020). Physics-informed recurrent neural networks for soft robots. \textit{arXiv preprint arXiv:2005.09333}.

\bibitem{lee2020model} Lee, K. C., \& Carlberg, K. T. (2020). Model reduction of dynamical systems on nonlinear manifolds using deep convolutional autoencoders. \textit{Journal of Computational Physics}, 404, 108973.

\bibitem{ljung1999system} Ljung, L. (1999). \textit{System Identification: Theory for the User} (2nd ed.). Prentice Hall.

\bibitem{vidyasagar1993nonlinear} Vidyasagar, M. (1993). \textit{Nonlinear Systems Analysis} (2nd ed.). Prentice-Hall.

\bibitem{hamilton1994time} Hamilton, J. D. (1994). \textit{Time Series Analysis.} Princeton University Press.

\bibitem{brockwell2016introduction} Brockwell, P. J., \& Davis, R. A. (2016). \textit{Introduction to Time Series and Forecasting} (3rd ed.). Springer.

\bibitem{pascanu2013difficulty} Pascanu, R., Mikolov, T., \& Bengio, Y. (2013). On the difficulty of training recurrent neural networks. In \textit{International Conference on Machine Learning (ICML).}

\bibitem{lim2021temporal} Lim, B., Ar{\i}k, S. {\"O}., Loeff, N., \& Pfister, T. (2021). Temporal fusion transformers for interpretable multi-horizon time series forecasting. \textit{Journal of Machine Learning Research}, 22(1), 1--41.

\bibitem{chen2021autoregressive} Chen, S., Rai, D., Thiel, R., \& Viswanathan, K. (2021). Autoregressive neural networks for the estimation of probability density functions. \textit{IEEE Transactions on Signal Processing}, 69, 1932--1946.

\bibitem{madaeni2017trajectory} Madaeni, S. H., Nouranizenouz, P., \& Zeadally, S. (2017). Trajectory-based sensor data compression for wireless sensor networks. \textit{IEEE Transactions on Wireless Communications}, 16(10), 6516--6529.

\bibitem{stock2001forecasting} Stock, J. H., \& Watson, M. W. (2001). Forecasting output and inflation: The role of asset prices. \textit{Journal of Economic Literature}, 39(3), 788--829.

\bibitem{lutkepohl2005new} L{\"u}tkepohl, H. (2005). \textit{New Introduction to Multiple Time Series Analysis.} Springer.

\bibitem{munkres2000topology} Munkres, J. R. (2000). \textit{Topology} (2nd ed.). Prentice Hall.

\bibitem{rufibach2018monotone} Rufibach, K. (2018). Monotone estimation of binomial probabilities. \textit{Computational Statistics}, 33(2), 531--552.

\bibitem{raissi2019physics} Raissi, M., Perdikaris, P., \& Karniadakis, G. E. (2019). Physics-informed neural networks: A deep learning framework for solving forward and inverse problems involving nonlinear partial differential equations. \textit{Journal of Computational Physics}, 378, 686--707.

\bibitem{li2021learning} Li, Z., Kovachki, N., Azizzadenesheli, K., Liu, B., Bhattacharya, K., Stuart, A., \& Anandkumar, A. (2021). Neural operator: Graph kernel network for partial differential equations. \textit{arXiv preprint arXiv:2003.03485}.

\bibitem{ragab2021attention} Ragab, M. G., Abdallah, Z., \& Azab, M. (2021). Channel attention based adversarial domain adaptation and generalization for unbiased learning. In \textit{IEEE Transactions on Industrial Informatics.}

\bibitem{jing2017tunable} Jing, L., Shen, Y., Dub{\v c}ek, T., Peurifoy, J., Sluran, S., Tegmark, M., \& Tegmark, L. (2017). Tunable efficient text mining with laplacian regularized linear regression. In \textit{Proceedings of the International Conference on Machine Learning.}

\end{thebibliography}

\end{document}